\section{Security Analysis}
\label{sec:security}

We analyze 6FLAS against the threat model of \cref{sec:threat}. Throughout, let
$\skey_\domid$ be the secret domain key; we model the keyed MAC as a secure
pseudorandom function (a computational assumption; see \cref{sec:formal}), and
recall that the Flow Label is immutable in transit.

\subsection{Deployment Model}
The analyses below assume all intermediate hops on a 6FLAS-protected
path implement \texttt{End.SEC}. In a mixed deployment where some transit nodes
implement only \texttt{End}, those nodes forward packets without performing security checks (since \texttt{End.SEC} is designed as a strict superset of \texttt{End} to maintain interoperability), meaning security guarantees hold only at the upgraded hops. This results in \emph{silent security degradation}: downstream 6FLAS nodes typically cannot detect that an intermediate plain \texttt{End} node skipped verification, because the packet format and tag remain intact through such hops. Mixed-deployment paths therefore provide best-effort, partial protection rather than end-to-end guarantees, and operators deploying 6FLAS should ensure complete coverage on security-critical path segments.

\subsection{Resilience to Known Attacks}
\label{sec:attacks}

We examine 6FLAS against four attack classes: forgery and tampering, replay, man-in-the-middle manipulation, and stateless denial of service.

\paragraph{Forgery and tampering (R1, R2).}
A packet is accepted only if the carried tag equals
$\mtag=\mac_{\skey_\domid}(\textsf{IPv6.src} \,\|\, \textsf{IPv6.dst} \,\|\, \fl\,\|\,\textsf{SID}\,\|\,\domid\,\|\,\seq)|_{32}$. An adversary without $\skey_\domid$ who
modifies any covered field must produce a matching tag.

\begin{proposition}[Forgery bound]
\label{prop:forgery}
Against an adversary making $q$ verification queries and lacking $\skey_\domid$,
the probability that any single forged or tampered packet is accepted is at most
$2^{-32}+\epsilon_{\prf}$, where $\epsilon_{\prf}$ is the PRF-distinguishing
advantage for the keyed MAC.
\end{proposition}

The $2^{-32}$ term is the chance of guessing the 32-bit truncated tag; a secure
MAC contributes only $\epsilon_{\prf}$. We prove this in \cref{sec:formal}.
Truncating to 32 bits trades tag length for header
compactness; a wider tag can be used when a stronger bound is required. Because the tag covers \textsf{SID}, an adversary also cannot
redirect a packet onto an unauthorized path without invalidating it (R4).

\paragraph{Replay (R3) --- probabilistic guarantee within window lifetime.}
Each accepted packet advances the per-$(\domid, \text{direction})$ sliding window $W_v[\domid, d]$. Since $\seq$ is a per-domain counter and the window is indexed by $(\domid, \text{direction})$, all flows from the same domain going in the same direction share the same replay window. This provides probabilistic anti-replay within the window lifetime: a replayed sequence number is rejected while it remains in the window, but once the window advances past it (due to newer packets from the same domain), the protocol itself no longer provides packet-by-packet replay detection. Periodic domain-key rotation invalidates old sequence spaces at a coarser granularity, providing long-term mitigation but not immediate per-packet rejection of old replays.

\begin{proposition}[Replay detection within window]
\label{prop:replay}
Let $R_{\domid,d}$ be the aggregate packet rate from domain $\domid$ in direction $d$ (i.e., the arrival rate into the replay bucket $W_v[\domid, d]$). A replayed packet with sequence $\seq_n$ is guaranteed rejected for at least
$\lfloor w / R_{\domid,d} \rfloor$ seconds after first acceptance, where $w=256$ is the window size.
\end{proposition}

For example, if $R_{\domid,d}=10^4$\,pkt/s, the window covers approximately $25$\,ms. Beyond the window lifetime, replayed packets with old sequence numbers are not detected by the window mechanism itself; they are invalidated by periodic domain-key rotation (typically hours to days), which resets the sequence space and makes all prior packets unverifiable.

The $w=256$-entry window tolerates up to $256$ out-of-order arrivals: a packet whose $\seq$ is ahead of the current base but within $w$
steps is accepted and its bit marked, so moderate reordering from multi-path
forwarding is handled without false rejection. For typical LEO packet rates and key rotation intervals, the combination of short-term window-based detection and long-term key-rotation-based invalidation covers the anti-replay requirement R3 at the rates and reordering levels the topology produces, while keeping the on-board state independent of the flow count.

\paragraph{Man-in-the-middle.}
An on-path adversary can observe and relay packets but cannot alter the covered
fields without breaking \cref{prop:forgery}, and cannot splice a valid tag from
one flow onto another because the tag binds the routing context (including IPv6 addresses and Flow Label) to the
SID and sequence number.

\paragraph{Denial of service.}
Stateful protocols are vulnerable to state-exhaustion DoS: an attacker opens
many bogus flows to fill the SA table. 6FLAS keeps no per-flow state, so this
vector is closed by construction. The first action in \cref{alg:verify} is a
domain-id check, and verification is a single MAC plus an $O(1)$ window probe,
so the per-packet work an attacker can induce is bounded and small. A Flow-Label
allow-list per domain (an optional deployment optimization) further lets a node reject traffic outside its assigned
label ranges before computing any MAC.

The four analyses show that 6FLAS satisfies R1--R4 under the
threat model: forgery is bounded by the MAC truncation (\cref{prop:forgery}), replay is
rejected within the window lifetime and invalidated long-term by key rotation (\cref{prop:replay}), man-in-the-middle manipulation is detected because the tag
binds routing context and path policy jointly, and state-exhaustion DoS is
eliminated by construction. The cost of each check is $O(1)$ in the flow count,
consistent with the feasibility condition of \cref{sec:formulation}.

\subsection{Comparative Security Properties}
\label{sec:sec-compare}

\cref{tab:secprop} compares the security properties 6FLAS provides against the
baselines. 6FLAS matches IPsec/DTLS on domain-origin authentication and integrity of the authenticated routing context,
provides probabilistic anti-replay within the window lifetime; it is the only scheme in the comparison that combines stateless verification
with topology-aware path access control. It deliberately omits payload
confidentiality, which can be layered on top via a composable encryption layer.

\begin{table}[t]
\centering
\caption{Security properties provided ($\bullet$\,=\,yes, $\circ$\,=\,partial, --\,=\,no).}
\label{tab:secprop}
\renewcommand{\arraystretch}{1.2}
\setlength{\tabcolsep}{4pt}
\begin{tabular}{lcccc}
\toprule
\textbf{Property} & \textbf{IPsec} & \textbf{DTLS} & \textbf{SRv6-Sec (HMAC)} & \textbf{6FLAS} \\
\midrule
Domain-origin auth.\ (R1) & $\bullet$ & $\bullet$ & $\bullet$ & $\bullet$ \\
Integrity (R2)          & $\bullet$ & $\bullet$ & $\bullet$ & $\bullet$ \\
Anti-replay (R3)        & $\bullet$ & $\bullet$ & $\circ^\dagger$   & $\circ^\ddagger$ \\
Path access control (R4)& $\circ$   & --        & $\circ$   & $\bullet$ \\
Stateless verification  & --        & --        & --        & $\bullet$ \\
Handover w/o re-key      & --        & --        & $\circ$   & $\bullet$ \\
Payload confidentiality & $\bullet$ & $\bullet$ & --        & -- \\
\bottomrule
\multicolumn{5}{l}{\footnotesize $^\dagger$ Assumes per-flow or per-policy state.} \\
\multicolumn{5}{l}{\footnotesize $^\ddagger$ Probabilistic within window; long-term via key rotation.}
\end{tabular}
\end{table}

\subsection{Formal Security Argument}
\label{sec:formal}

The authentication guarantee is made precise by reducing it to the
unforgeability of the keyed MAC.

\begin{proof}[Security argument for \cref{prop:forgery}]
We model keyed SipHash-2-4 as a secure pseudorandom function; this is a
computational assumption --- PRF security is not a proven property of SipHash-2-4
but is the standard working assumption for keyed SipHash in protocol
analysis~\cite{siphash}. Under this assumption, $\mac_{\skey}|_{32}$ is a
secure truncated MAC by standard length-extension-free truncation arguments.
Let $\textsf{MAC}_{\skey}$ be the keyed MAC of \cref{eq:tag} and let
$\textsf{MAC}_{\skey}|_{32}$ be its $32$-bit truncation. Consider an adversary
$\mathcal{A}$ that, after observing honestly generated packets, submits a packet
whose covered message $M=(\textsf{IPv6.src}\,\|\,\textsf{IPv6.dst}\,\|\,\fl\,\|\,\textsf{SID}\,\|\,\domid\,\|\,\seq)$ was never
produced by an honest source. (If $M$ equals an honest message, the packet is a
replay and is rejected by the sliding window of \cref{alg:verify} with
high probability within the window lifetime (see \cref{prop:replay}), so we may assume $M$ is
fresh.) Acceptance requires the carried tag to equal
$\textsf{MAC}_{\skey}(M)|_{32}$.

Build a forger $\mathcal{B}$ against $\textsf{MAC}_{\skey}$: $\mathcal{B}$
forwards $\mathcal{A}$'s observation requests to its own tagging oracle to
generate honest packets, and outputs $\mathcal{A}$'s fresh pair
$(M,\mtag)$ as a forgery. Because $M$ was never queried, this is a valid forgery
attempt. If $\textsf{MAC}_{\skey}$ is a secure PRF with distinguishing advantage
$\epsilon_{\prf}$, the probability that $\mtag$ matches the true full tag is at
most $\epsilon_{\prf}$ beyond random guessing. With $32$-bit truncation, an
adversary that simply guesses succeeds with probability $2^{-32}$. Combining,
the per-packet acceptance probability of a fresh forgery is at most
$2^{-32}+\epsilon_{\prf}$. By a union bound over $q$ independent forgery
attempts, the overall bound is $q\,(2^{-32}+\epsilon_{\prf})$.
\end{proof}

The reduction shows the structure of the guarantee: a successful network
adversary $\mathcal{A}$ that produces an accepted packet it never observed from
an honest source is exactly a MAC forger, because the covered fields it submits
differ from any honestly produced message (otherwise the packet is a replay,
caught by the window). Hence 6FLAS's authentication security reduces to the
unforgeability of its keyed MAC, with the $32$-bit truncation contributing the
$2^{-32}$ term of \cref{prop:forgery}.
